\documentclass[12pt]{article}
\usepackage{graphicx} % Required for inserting images
\usepackage{authblk}      % Optional: better author formatting
\usepackage{setspace}     % Optional: spacing control
\usepackage{fancyhdr}
\usepackage{amsmath}
\usepackage{caption}
\usepackage{pdflscape}
\usepackage{graphicx}
\usepackage{booktabs}
\usepackage{adjustbox}
\usepackage{url}
\usepackage[utf8]{inputenc}
\usepackage{csquotes}
\usepackage{natbib}
\usepackage{pdfpages}
\usepackage[raggedright]{ragged2e}
\usepackage[hidelinks]{hyperref}
\usepackage[letterpaper,margin=1in]{geometry}
\bibliographystyle{apsr}   % or plainnat, chicago, apa, etc.
\setcitestyle{authoryear,open={(},close={)},aysep={ }}

\title{We Don't Do That: State Identity and Nuclear Non-Use\footnote{Presented to the Columbia University Department of Political Science for satisfactory completion of the requirements to graduate with departmental honors.} \footnote{I thank my thesis advisor, John Marshall, honors seminar proctor, Shigeo Hirano, and graduate teaching assistant, Indira Tirumala, for their feedback in developing all aspects of this paper. I am grateful to the Department of Political Science, Undergraduate Research and Fellowships, and the Columbia Experimental Laboratory for the Social Sciences for providing funding to support this project. Lastly, I thank my friends—Conrad Hutchins, John Garby, Nicole Wong, Oscar Manuel Landa Samano, Erin Stoeckig, and many, many others—for entertaining my frequent digressions about nuclear weapons and Jawaharlal Nehru.}} 
\author{Wyatt King}
\date{Draft Sections: Introduction, Literature Review, Theory, Case Selection, Research Methods, February 27}


\begin{document}

\maketitle
\begin{abstract} 
\noindent
This is where your abstract goes. It should briefly summarize the research question,
methods, and findings.
\end{abstract}

\newpage


\section{Introduction}
In international relations scholarship, political scientists have often taken after microeconomic theory in assuming that states are rational and selfish \citep{waltz_man_2001}. Neorealist scholars presuppose that states are strictly security seeking, unmoved by considerations of fairness or altruism in an international system predicated on anarchy and self-help. Since the “constructivist turn,” however, scholars have challenged the descriptive accuracy of assuming states to be egoistic and the international system anarchic \citep{wendt_social_1999, finnemore_international_1998, wendt_anarchy_1992}. Instead, individual leaders, domestic constituencies, and states as units in and of themselves can be constrained by norms of how they are meant to behave. 

Crucial to how norms are theorized to affect state behavior is identity. If an action is so morally reprehensible as to contradict a state's perception of itself, then it will refrain from whatever that action is. Because identity and norms are difficult to measure numerically, constructivist scholars have often resorted to process-tracing and other qualitative means to substantiate their argument that states' behavior is mediated by ideational factors \citep{peez_contributions_2022}. Although these approaches have been useful in illustrating how identity affects elites' attitudes toward nuclear use \citep{tannenwald_nuclear_2007}, collateral damage \citep{crawford_accountability_2013}, and other norm-breaking behavior, it has limited the scope of constructivist epistemology. The lack of quantitatively-driven, positivist study is striking given that experimental research has shown priming national identity can alter political behavior, such as by reducing affective polarization \citep{levendusky_americans_2018} or changing attitudes toward immigration \citep{wojcieszak_social_2018}. In this paper, I seek to redress this methodological gap by fielding a survey experiment in which Indian respondents are exposed to identity-based arguments opposing a hypothetical counterforce attack against Pakistan. 

My focus on nuclear weapons follows from scholarship that has characterized non-use as among the strongest norms in the international system. Describing use as repulsive and “taboo,” Nina \citet{tannenwald_nuclear_2007} contends that non-use is reflexive. Even in dire circumstances, a shibboleth is attached to nuclear weapons, effectively eliminating states' motivation to use them. Realists have criticized Tannenwald's argument, saying that if a norm against non-use exists at all, it is epiphenomenal—the product of egoistic states acting out of self-interest \citep{sagan_realist_2004}. To substantiate their argument that no taboo exists, Scott Sagan and Benjamin Valentino have conducted survey experiments revealing that Western public opinion is far more receptive to the use of nuclear weapons than previously thought, with compelling majorities approving of their use in circumstances where they achieve a compelling strategic objective \citep{press_atomic_2013, sagan_revisiting_2017, dill_kettles_2022}.

By focusing on India, I test the salience of identity-based arguments in the circumstance where they are most likely to have an effect. Unlike in Western, nuclear armed states, political elites in India have historically framed non-violence as morally obligatory and a touchstone for its foreign policy \citep{perkovich_indias_2001, hymans_psychology_2006, mehta_still_2009, bajpai_indian_2014}. Following World War II, India was among the staunchest proponents of global disarmament and continued to be so even after it tested nuclear weapons. Today, although political elites raise doubts that India will abide by its no first use policy, they still explicitly assert that first use is antithetical to Indian culture and civilization.

Aside from entertaining how identity affects attitudes toward nuclear weapons use, I consider how the institutionalization of non-use can do the same. Whereas previous studies have focused on how international law can restrict support for nuclear use \citep{carpenter_stopping_2020, sagan_atomic_2024}, I consider how no first use policies (NFUs)  do the same. Whereas international law promotes shared norms, NFUs reinforce perceptions of state's personal identities as being incompatible with nuclear use. This line of argument presents a novel modification to scholarly understanding of NFUs' credibility. Specifically, it raises the possibility that NFUs make nuclear use more politically costly by reifying an image of state identity as being incompatible with nuclear use.

Ultimately, the purpose of this thesis is twofold. Narrowly, it contributes to the upswell of survey experiments from the past ten years, in which scholars have sought to understand what factors impact public support for the use of nuclear weapons. It does so by interrogating the effects that national identity and institutionalization can have, and whether those effects are significant in the manner that scholars of the taboo suggest. More broadly, it seeks to evaluate claims in constructivist literature that norms can have a “constitutive effect,” in which people's foreign policy preferences are impacted by their impression of state identity.

[DISCUSSION OF RESULTS]


\section{Literature Review \& Theory}

In international relations scholarship, there are two broad theoretical orientations that account for why states behave in the manner that they do: Rationalism and Constructivism \citep{katzenstein_international_1998}. Though not wholly incompatible, they differ in their understanding of states' interests and the power of norms in the international system \citep{fearon_skeptical_2002}. As part of these distinct interpretations of how states behave, rationalists argue that all states are fundamentally like one another, doing whatever advances their material interests. Constructivists counter that states have distinct identities which inform how they behave toward others.

Per rationalist scholarship, states exclusively seek to improve their overall material condition \citep{waltz_man_2001, keohane_after_2005}. They may want either security, power, or wealth, but regardless of what they are pursuing, states are rational actors optimizing in the name of that interest. They are egoistic, concerned with other states only insofar as they may be an impediment to or an instrument for the realization of a strategic objective. If norms exist at all in the international system, they are epiphenomenal \citep{sagan_realist_2004}. That is to say, they are not rooted in moral reasoning but stem from states acting in a manner consistent with whatever they believe to be in their material interests. In this formulation, norms do not actually regulate state behavior but are the illusory product of states' behaving rationally.

Because rationalists assume that states desire comparable material objectives, there is no role for identity in shaping policy. Rather, all states are like one another in their disposition. For realists, the international system is like Hobbes' state of nature. Unconstrained by an overriding legal authority, states must protect themselves and their material interests using whatever means are available to them. Should a norm-breaking action, like the use of nuclear weapons or the deliberate killing of civilians, suit a states' perceived short- and long-term interests, they will follow that course of action. Neoliberals are more optimistic than realists that states can effectively cooperate to overcome collective action problems, such as non-proliferation, but they operate from the same underlying assumptions about rationality and egoism. Instead of being moved by principles of what they ought to do, all states are utility-maximizers doing whatever satisfies their material interests.

Constructivism problematizes rationalists' dominant assumption that states seek to guarantee their security or otherwise improve their material condition \citep{wendt_social_1999}. Instead of treating the nature of the international system as immutable, constructivists treat the international system as being socially constructed \citep{finnemore_international_1998}. If states understand security as zero-sum or wealth as their ultimate objective, that is because they have been cultured into perceiving their interests in this manner. Because the constitution of the international system is a social phenomenon, constructivists argue that norms can influence states' behavior. As commonly understood, norms are “a standard of appropriate
behavior for actors with a given identity,” \citep[891]{finnemore_international_1998}. If an action violates an actor's self-perceived identity, then they will refrain from whatever that action is. In simplified terms, arguments about state identity and morality amount to saying “\textit{we} don't do \textit{that}.” Because of some deeply rooted understanding of self-hood, certain behaviors are considered anathema to a state's interests \citep{wendt_social_1999}.

In constructivist scholarship, identities are inter-subjective and shared between like-minded states, producing a similar pattern of behavior among states who claim that identity. States are incentivized to comply with norms because if they do not, they will experience embarrassment, shame, anxiety, or guilt among their peers \citep{fearon_what_1999}. As well as being social, norms can also be personal. States may be expected to act in a given way if they are “Christian,” “Marxist,” or “post-colonial,” but these titles convey little about the cultural, economic, and historical conditions that make them unique relative to other states with whom they share group identities. These personal identities create another opportunity for norms to influence state behavior. Even if a given action does not violate a norm shared by an over-arching group, it can violate an aspect of one state's identity from which it derives pride, a feeling of righteousness, or a similarly salient affect. That being said, because personally-held identities are exclusive to individual states, they are theorized to be more malleable than group identities, subject to reinterpretation if and when a norm-breaking action suits a state's material interests. 

In debates between rationalists and constructivists, no single security issue has prompted as much discourse as nuclear weapons use \citep{paul_nuclear_1995, sagan_realist_2004, tannenwald_nuclear_2007, paul_nuclear_1995, smetana_forum_2021}. According to balance of power realists, nuclear non-use is attributable to the achievement of mutually assured destruction during the Cold War. Because no state can use atomic weapons without fearing annihilating retaliation, it was in no one's material interests to use nuclear weapons. For however much MAD explains relations between nuclear-armed opponents, it does not explain the lack of nuclear use in conflicts between nuclear-armed and non-nuclear states. Why, for example, did the U.K. not use nuclear weapons during the Falklands War, or the U.S. during the first Gulf War? In each instance, there is reason to think that nuclear weapons would have been useful on the battlefield, yet even discussion of their use was “taboo” \citep{paul_nuclear_1995, tannenwald_nuclear_2007}. 

Constructivists attribute the \textit{de facto} prohibition on the first use of nuclear weapons to a widespread, moral revulsion attached to them. A taboo transcends a norm because breaching it is unthinkable and constitutes crossing a bright red line. Breaking the nuclear taboo is not only something that states avoid doing, as they might eschew deliberately killing civilians, but it is even unimaginable. \citet{tannenwald_nuclear_2007} argues that the taboo developed, in part, because states began to frame their identity as being antithetical to nuclear weapons use. During the first Gulf War, for example, Chairman of the Joint Chiefs of Staff Colin Powell called using nuclear weapons against Iraq “un-American” and tore up a war plan describing the U.S.'s nuclear options, saying that he did not want to “let the Genie out of the bottle.” When asked by the press if the U.S. would consider nuclear use, White House Chief of Staff John Sununu said it was something, “we just don't do.”


Today, the most lively debate between rationalists and constructivists over nuclear non-use takes the form of public opinion research \citep{smetana_forum_2021}. \citet{tannenwald_nuclear_2007} argues that although political elites ultimately decide whether nuclear weapons should be used, staunch public opposition can limit the probability that they are. To substantiate her claim that the public is, in fact, averse to nuclear use, Tannenwald points to surveys showing declining support for the U.S.'s bombings of Hiroshima and Nagasaki since the 1940s.\footnote{See \citet{sagan_revisiting_2017} for a condensed overview of this polling.} Moreover, polling conducted concurrently with various conflicts suggest that there was broad public opposition to using atomic weapons. Before the first Gulf War, for example, more than 75\% of respondents opposed using nuclear weapons should the war have become a stalemate, and a slim majority opposed their use even if Iraq attacked American troops with chemical or biological weapons (CNN/Time Magazine 1990). 

Over the last ten years, skeptics have challenged whether there is a widespread, public abhorrence of nuclear weapons. Focusing on methodological innovations instead of theoretical ones, new research has employed survey experiments to evaluate the extent of public support for a first strike nuclear attack in various hypothetical scenarios. \citet{sagan_revisiting_2017} notably simulated what they argue is a modern-day equivalent of the U.S.'s bombings of Japan. Contending that polling Americans about their retrospective preferences for nuclear use during World War II does not adequately prime them to think about the trade-offs of restraint, they presented respondents with a scenario in which the U.S. dropped an atomic bomb on Iran's second largest city to end a ground war that would otherwise kill 20,000 American troops. Despite respondents being told that such an attack would kill two million civilians, 60\% of people approved of it.

Following the publication of these findings, similar studies have since been released affirming that large swaths—if not outright majorities—of the Western public would approve of the first-use of nuclear weapons in different hypothetical scenarios \citep{aronow_note_2019,haworth_what_2019, dill_kettles_2022,schwartz_when_2024}. As these studies have proliferated, the scope of debate has expanded beyond whether there is a broad based public opposition to nuclear use to what factors can attenuate support for it. In answering this question, scholars have ignored identity, focusing narrowly on three factors: The military and strategic utility of a nuclear attack, states' legal obligations under international law, and the number of civilian casualties incurred by an attack.

\textit{Military and Strategic Utility} Among the first considerations that Sagan and Valentino entertained in their on-going effort to ascertain what factors affect public preferences for nuclear use was the military and strategic utility of an attack. Together with Daryl Press, they argue that the public overwhelmingly—though not exclusively—follows a utilitarian logic of consequences in weighing whether to support nuclear use \citep{press_atomic_2013}. According to the logic of consequences, people's support for nuclear use is predicated on their perceptions of the combined short- and long-term advantages of nuclear weapons as compared to conventional ones. By the military utility of an attack, Press, Sagan, and Valentino refer to the immediate battlefield advantages conferred by atomic weapons use. Should a nuclear bomb destroy an enemy or their infrastructure more substantially, with greater certainty, and while better mitigating domestic casualties than conventional arms, then people will be more inclined to support their use.

By the strategic utility of an attack, the authors refer to the long-run disadvantages tied to setting a new precedent for nuclear use. As Paul (2009) argues, nuclear non-use can be understood as an iterated prisoner's dilemma. On any one battlefield, it may be in a state's interests to defect and use nuclear weapons, but doing so incentivizes an adversary to do the same in a future conflict. As a result, nuclear use may not be desirable because of how it transforms opponents' utility functions during future crises. According to this formulation, the apparent norm against nuclear use is epiphenomenal. Rather than stemming from moral reasoning, it is the product of coordination, either implicit or explicit, between nuclear-armed states to achieve a mutually beneficial outcome.

Survey experiments have consistently shown that military and strategic calculations exert a powerful influence on the public's willingness to support nuclear weapons use. When a state has a nuclear-armed opponent, the public is far less willing to use nuclear weapons than in asymmetric dyads \citep{haworth_what_2019,allison_under_2022, smetana_moscow_2022}. Holding the opponent constant, support for nuclear use increases dramatically with the chances of successful attack relative to conventional weapons \citep{press_atomic_2013}. To test whether the strategic utility of an attack impacts public preferences for nuclear use, \citet{sagan_atomic_2024} replicated their\textit{ Hiroshima in Iran} scenario, but informed respondents that members of the Joint Chiefs of Staff worried that using weapons would set a bad precedent that could ultimately harm American security. In this treatment group, approval of the attack declined by nearly fifteen percentage points.

\textit{International Law}
\noindent
Beyond the military and strategic utility of a nuclear attack, scholars have considered how international law may limit states' willingness to use nuclear weapons. According to \citet{tannenwald_nuclear_2007}, the institutionalization of nuclear non-use through international agreements has concretized the taboo against atomic weapons. Tannenwald interprets institutionalization broadly, saying that it spans not only international law but also the Nuclear Non-Proliferation Treaty (NPT), bilateral anti-ballistic missile agreements, and declared No-First Use policies (NFUs). In survey experiments, however, scholars have primarily focused on international law.

According to the logic of consequences, international law may impact public opinion if people believe that, by breaking it, their state will face material repercussions for their actions. Because international law lacks an enforcement mechanism, though, \citet{sagan_does_2020} argue that it likely has little effect on public attitudes toward nuclear use. Constructivists, however, have argued that international law may have an effect on state behavior if states follow a logic of appropriateness rather than one of consequence. According to the former, states are rule followers who, though not obligated to behave in a given manner, do so anyways to validate a social identity that they have constructed (Fearon and Wendt 2002). According to this thinking, states are not strictly security seekers as neorealists suggest. Rather, they are bound by an informal set of ideas and norms for how states similar to them are expected to behave. 

Outside of nuclear use, surveys have found that the American public is reluctant to support actions knowing that they violate international law \citep{chilton_laws_2014, kreps_flying_2014}. \citet{carpenter_stopping_2020}, replicating a scenario from \textit{Hiroshima in Iran} in which the U.S. uses conventional weapons in a saturation bombing, find that respondents' support for the bombing declines after they are informed that it violates international law. \citet{sagan_does_2020} question the relevance of their finding, however, by pointing out that the effect size of informing people about international law is relatively small and easily reversible. Replicating Carpenter and Montgomery's research, they found that while a minority of respondents approved of the saturation bombing after being informed of it violating international law, only slightly less than 50\% actually did. In a follow up looking at nuclear use, \citet{sagan_atomic_2024} note that if there was disagreement among the Joint Chiefs of Staff regarding the legality of a strike, the initial effect of telling people about international law washed out.


\textit{Civilian Casualties and Moral Frameworks}
\noindent
Scholars have considered how the number of civilian casualties may change the extent of support for nuclear use. If it is true that states have internalized a sensitivity to excessive civilian casualties, owing in part to the legacy of aerial bombings during World War II \citep{pinker_better_2011, thomas_ethics_2014}, then it should follow that people would be opposed to a nuclear attack that would cause large-scale, unnecessary violence. On theoretical grounds, \citet{downes_desperate_2006} and \citet{sagan_revisiting_2017} dispute this characterization by arguing that if such a norm exists at all, it is not strong enough to withstand the domestic political pressure that would arise from a costly war. Downes even concludes that democratically-elected governments are more likely to commit violence against civilians because they may be sensitive to the costs of conflict and, thus, more likely to resort to desperate means. Sagan and Valentino, focusing more narrowly on public opinion, make a similar point, arguing that voters are willing to tolerate large-scale violence perpetrated against civilians in the event substantial troop casualties would otherwise occur.

To test whether civilian immunity exerts a constraining influence on nuclear use, \citet{sagan_revisiting_2017} manipulated the number of civilian casualties incurred by a first-strike attack, holding the number of domestic casualties in a would-be ground war constant at twenty thousand. Despite ratcheting the number of expected civilian deaths from one hundred thousand to two million, public approval of the attack decreased by only a fraction of a percentage point. Meanwhile, preference for the nuclear attack over continuing the ground war declined by eight percentage points, a margin that was not statistically significant in their sample.

In response to this finding, \citet{rathbun_greater_2019} argue that there are a diversity of moral frameworks that people may use in evaluating the merits of nuclear use and not strictly a logic of consequence or appropriateness as Press, Sagan, and Valentino initially argued. Using the typology of moral foundations developed by Graham et al. (2011), they found that people's support for nuclear use changes depending on whether they adopt liberal, non-liberal, or retributive moral frameworks. Based on this typology, Rathbun and Stein subset  respondents on whether they had “individualizing foundations,” a sensitivity to others' suffering and a concern for equality and fairness, or “binding foundations,” which emphasize deference to authority, in-group loyalty, and the subordination of individual rights to the community. Among those who had stronger individualizing beliefs, support for nuclear decreased substantially as civilian causalities rose. That being said, few respondents in the total population preferred using nuclear weapons to conventional ones as the number of civilian casualties approached one million, in contrast to Sagan and Valentino's earlier finding.

Cumulatively, research on the public's aversion to nuclear use has shown that its attitudes are not deontological. Rather than there being a taboo, the public is supportive of nuclear use to prevent domestic casualties and achieve a compelling military objective. Yet, even as extant surveys have revealed that the public does not possess a categorical aversion to nuclear use, there is reason to believe that a norm exists it in some capacity. Holding all else constant, American respondents vastly prefer the use of conventional weapons to nuclear ones \citep{press_atomic_2013}. Moreover, even as existing studies have created “tough tests” for the taboo, in which atomic weapons use would be most compelling, nearly half of the public still disapprove of it \citep{sagan_revisiting_2017}.

Although survey experiments have helped clarify what factors impact American attitudes toward nuclear use, they have yet to test the underlying tension between constructivist and rationalist approaches: that states' identities are socially constructed and can influence their understanding of what is appropriate conduct. If it is the case that Americans are receptive to strategic nuclear use, impervious to the number of civilian causalities, and apathetic about the legality of first-use, then that may not mean all states are similarly materialistic. Rather, it could imply that Americans do not configure their understanding of state identity as being antithetical to violating the nuclear taboo, international law, or civilian immunity.

To be sure, \textit{ex ante}, there is reason to believe that the U.S. is unique in its understanding of identity and international norms. That the U.S. is the only state to have ever used nuclear weapons outside of testing may limit how much the public understands doing so again as “un-American.” Likewise, even though the U.S. has not used atomic weapons since 1945, arms-racing during the Cold War and near-uses, such as during the Cuban Missile Crisis, may acclimate people to the possibility that nuclear weapons could and should be used when necessary. In a survey experiment of the four Western, nuclear armed states, \citet{dill_kettles_2022} found that, although no group exhibited an overwhelming aversion to nuclear use, there were statistically significant differences in their degree of support. British survey takers were less likely to approve of a limited nuclear strike than American or French ones. Israeli respondents, meanwhile, were more hawkish, both more inclined to approve of the attack and say that it were ethical than all other groups. Though this research does not necessarily mean that identity is the mediating variable causing differences in support for nuclear use, it raises the possibility that it could be. As a result, further empirical study is required to understand the breadth of how identity can impact public attitudes toward nuclear use.

In this paper, I propose advancing the experimental study of identity and nuclear use in three ways: by sampling more states in further research, testing how identity-based arguments alter support for nuclear use, and measuring how the institutionalization of nuclear non-use reinforces identity-based appeals to mitigate the public's overall approval of nuclear use.

On the first count, although Dill, Sagan, and Valentino made an initial effort to understand cross-national differences in support for nuclear use, constructivists would not expect \textit{a priori} that attitudes toward it would differ meaningfully across Western, nuclear armed states. If identities are inter-subjective and shared among states, then historically close allies with similar styles of government should behave like one another. Whether it is because they have a common understanding of what a “liberal” or “civilized” state does, the overall compatibility of their identities may explain why, despite there being statistically significant differences in their degree of support for nuclear use, they were not substantively different from one another.

To evaluate support for nuclear use in a different state with a distinct identity, I propose focusing on the one outstanding nuclear-armed democracy: India. Focusing on India is useful theoretically for two reasons. First, because it is has an elected government, there is a compelling mechanism through which public preferences for nuclear use can impact elites' actions, unlike in China, Russia, North Korea, or Pakistan. In keeping with Tannenwald's (2007) initial theorizing on the taboo, I assume that election-seeking politicians will avoid nuclear use if they fear reprimand from their domestic audience. Second, as I will elaborate on in the following section, India has a history of political elites framing non-violence as morally obligatory and a touch stone for its foreign policy. As it manifests in nuclear doctrine, India has been among the most outspoken proponents of global disarmament. Even after it tested nuclear weapons, successive governments advanced plans for disarmament in international forums, often framing their support for disarmament in the terms of Gandhianism. Consequently, among nuclear armed states, the norm against nuclear use has historically been most strongly entrenched in India.

If India is the most likely case in which a norm against nuclear use may be observed, then the degree of absolute support is important. Should opposition to nuclear use not be overwhelming, then yet further doubt would be cast on the argument that there is a taboo against nuclear weapons use in any case. To formalize the expectation that there is a widespread, categorical aversion to the use of nuclear weapons in public opinion, I test how strongly the public would support nuclear use in the circumstance most favorable for their use. Here, I take that circumstance to be a successful counterforce attack against Pakistan that eliminates its strategic nuclear weapons. This scenario constitutes the toughest, reasonable test for the nuclear taboo in India because should it eliminate Pakistan's strategic nuclear weapons, then it could use its conventionally stronger military to take disputed territory in Kashmir without fear of excessive civilian casualties. As such, my first hypothesis can be written as follows: 

\begin{quote}
\textbf{H1}: Opposition to a first strike counterforce attack against Pakistan should be high among the Indian public.
\end{quote}

Regardless of if opposition to nuclear use against Pakistan may be high, it is worth considering what factors could possibly expand support for non-use. Following from constructivist theory, perceptions of state identity could be one such factor. If it is the case that states' conduct is limited by their conception of self, then it follows that invoking a given understanding of state identity will cause people to oppose behaviors incompatible with that identity. As a result, I hypothesize that:

\begin{quote}
\textbf{H2:} Public support for the use of nuclear weapons will decrease as people are exposed to arguments that nuclear weapons use is antithetical to India's state identity.
\end{quote}

Beyond the role that identity plays in shaping state behavior, constructivist scholarship argues that the institutionalization of norms reifies them, potentially leading to their widespread internalization (Finnemore and Sikkink 1998). As mentioned, a handful of survey experiments have looked into how priming survey respondents about international law impacts their support for nuclear use \citep{carpenter_stopping_2020, sagan_does_2020, sagan_atomic_2024}.  Focusing on international law is worthwhile and contributes to scholarly debates outside of nuclear weapons strategy regarding its effectiveness in regulating state behavior; however, pursuant Tannenwald's initial theorizing, it is worth considering how other types of institutions affect support for nuclear use. 

I propose focusing on No First Use policies, specifically. NFUs are unique among types of institutionalized non-use because they are explicit and unilateral. Compared to international law, which does not explicitly ban the use of small-yield tactical nuclear weapons or similarly proportionate nuclear attacks, NFUs do, amounting to an almost absolute prohibition on the first-use of nuclear weapons.\footnote{ NFUs may permit nuclear retaliation against a chemical and biological weapons attack. They may also permit a preemptive first-strike should an adversary credibly threaten an imminent nuclear attack.} As it relates to identity, NFUs are unique in that they are unilaterally declared. Rather than being the product of international agreement, NFUs are exclusively held by one state, meaning that rather than speaking to an inter-subjective, shared identity, they evoke a personal identity that states claim. In the case of India, its NFU overlaps with the commonly made nationalist claim that it is a unique moral authority in the international system and a “responsible” nuclear power, reinforcing the salience of arguments that Indian identity is incompatible with nuclear use. As such, I hypothesize that:

\begin{quote}
\textbf{H3:} Informing people about India's No First Use policy will decrease support for a first-strike nuclear attack.
\end{quote}

\begin{quote}
\textbf{H4:} Informing individuals about India's NFU in combination with exposure to pacifistic national identity arguments will reduce support for the use of nuclear weapons more than either treatment alone.
\end{quote}

Aside from directly testing rationalist and constructivist theories, these hypotheses bear implications for audience cost theory by answering whether declaring an NFU or invoking a pacifistic state identity cheap talk. In democracies, should either spark backlash against a government, then they are not. Rather, to borrow terminology from James \citet{fearon_domestic_1994}, they confer audience costs. These audience costs, however, are of a different variety than those that Fearon discusses. Rather than audience costs being paid by leaders who back down from a threat, they are paid by a leader who contradicts a state's identity as non-violent, responsible, or similar. Previous research on “consistency costs” suggests that public opinion punishes leaders for entering into a war that they previously promised to avoid \citep{levy_backing_2015}. Here, however, I argue that domestic public opinion is not exclusively sensitive to an individual leader's consistency but to the broader social construction of state identity. As a result, even if an individual leader were to disingenuously evoke an understanding of state identity as being antithetical to violence, they may still 

Taken together, hypotheses two through four represent a  modification to audience cost literature. As theorized by James \citet{fearon_domestic_1994}, audience costs are the cumulative domestic reprisal that leaders face for reneging on threats against a non-compliant adversary. Crucially, audience costs according to Fearon are paid by leaders who pledged to be violent in the face of a challenger but were not in actuality. As \citet{levy_backing_2015} argue, audience costs are not exclusively paid by actors who are insufficiently aggressive. Rather, they are also paid for by leaders who pledge to be pacifistic in the face of a threat, but subsequently support going to war. Explaining different types of domestic reprisal, Fearon and Levy et al. provide distinct rationales for why states pay audience costs. According to Fearon, a leader who is insufficiently violent will face audience costs because in making a threat against another state, they “engage the national honor” only to impeach it upon backing down. Though Fearon does not provide a definition of what he means by “honor,” in everyday parlance, it refers to the respect that someone derives from their merit, reputation, or character. In his discussion of honor, Fearon does not entertain sociological literature, but honor is widely understood to be socially constructed, with honor cultures varying both across and within states. The social construction of honor presents a potential caveat to Fearon's causal reasoning because should one person or state derive pride from being responsible, then they not pay audience costs for backing down from what are seen as irresponsible threats.

To explain why leaders also pay costs for reneging on commitments to stay out of war, Levy et al. argue that the public values consistency in leadership. Should a leader flip flop between policy preferences, it makes them appear indecisive or non-committal. My hypotheses both drawn from Fearon's theory and enhance Levy et al.'s argument by suggesting that identity and perceptions of honor may also be a mechanism through which leaders are punished for supporting violent foreign policies. If a state derives honor from being restraint and responsible, then 

\section{Case Selection}
To test whether identity can alter support for nuclear first use, I field a survey experiment modeled off of those created by Sagan and Valentino. In the experiment, respondents read a hypothetical scenario describing an Indian counterforce attack against Pakistan. I conduct the experiment in India because aside from it being the only democracy with a declared No First Use policy, political elites have framed India as a uniquely responsible and moral power. As it relates to nuclear use, diplomats and elected politicians have framed atomic weapons as being incompatible with both India's personal identity and an inter-subjective identity shared by post-colonial states. By the former, I refer to how leaders of the Indian nationalist movement adapted non-violence to Indian foreign policy at least rhetorically. By the latter, I refer to how India has defined itself in contrast to the U.S., U.K., and other nuclear weapons states. Whereas others have often stymied efforts to achieve global disarmament, India has embraced disarmament, criticizing others for promoting “nuclear apartheid” through the Non-Proliferation Treaty and the Comprehensive Test Ban Treaty.

The incorporation of non-violence into Indian national identity stems from the nationalist movement, in which political elites argued that non-violence was both an effective anti-imperial strategy and a moral obligation. None did so as clearly as Mahatma Gandhi. According to Gandhi, people were spiritually required to uphold \textit{ahimsa}, a religiously-inflected understanding of non-violence shared by Hindus, Jains, and Buddhists that literally means “to not kill or injure” \citep{young_hinduism_2004}. Inspired by an ascetic interpretation of Hinduism, Gandhi's understanding of \textit{ahimsa} was deontological and made no meaningful distinction between offensive and defensive wars. Following from his categorical aversion to violence, Gandhi's opposition to nuclear use was staunch and unequivocal. Reflecting on the American bombings of Hiroshima and Nagasaki, he said that, “Unless now the world adopts non-violence, it will spell certain suicide for mankind” \citep[12]{mitra_moral_1987}. When challenged that nuclear weapons could be an instrument to achieving peace by creating stable deterrence, Gandhi was incisive. He compared nuclear weapons' deterrent effect to a “a man glutting himself with dainties” \citep[278]{gandhi_my_2025}. States will refrain from violence “only to return with redoubled zeal after the effect of nausea.”

Following Gandhi's death, his peers in the nationalist movement adapted non-violence to Indian foreign policy, at least rhetorically. In 1954, after the U.S. and U.S.S.R. tested hydrogen bombs, Prime Minister Jawaharlal Nehru became the first world leader to call for disarmament. During negotiations over what would become the Nuclear Non-Proliferation Treaty, India sought to give teeth to these demands. Criticizing nuclear-armed powers for their hypocrisy in wanting to stop proliferation while allowing themselves to maintain and enlarge their nuclear arsenals, India refused to sign a multilateral non-proliferation treaty unless nuclear weapons states were required to disarm. Otherwise, the NPT amounted to what Indian policymakers called “nuclear apartheid,” in which wealthy states would secure their access to nuclear weapons and protect a crucial strategic advantage over poorer, primarily post-colonial states. When the U.S., Soviet Union, and others rejected calls to disarm, India refused to join the NPT.

To be sure, for however much India's opposition to the NPT stemmed from sincere moral feeling, it was also cut with an understanding of strategic necessity \citep{perkovich_indias_2001}. In 1964, China conducted its first nuclear tests and proceeded to build up its arsenal. Due to long-standing border disputes between the two countries, which culminated in the 1962 Sino-Indian War, China's possession of nuclear weapons risked tilting the balance of power in the Himalayas in its favor. By not signing the NPT, India guaranteed itself the option to later balance against China by testing nuclear weapons—as it would in 1974. India's balance between morally-driven foreign policy and \textit{realpolitik} was perhaps best summarized by Prime Minister Indira Gandhi, who in a speech before parliament in 1968 said that “we shall be guided by our self-enlightenment and the considerations of national security.”

Even as India pursued nuclear weapons, it continued to frame itself as a non-violent and responsible power. The codename for India's nuclear test was “smiling Buddha,” and in announcing its success, Indira Gandhi deemed it a “peaceful nuclear explosion.” Aside from naming conventions that could amount to cheap talk, India did not press its nuclear advantage over Pakistan after testing. It did not invest in ballistic missile systems or submarines that would make its deterrent more credible, and it did not test nuclear weapons again until more than two decades later. Moreover, India continued its advocacy for global disarmament. Indira Gandhi coordinated a group of five other states to advocate for disarmament, and in 1988, her son, Prime Minister Rajiv Gandhi, proposed a plan at a special U.N. summit on disarmament that would eliminate nuclear weapons globally by 2010. 

Compared to Congress-led governments of the mid- to late-twentieth century, the Bharatiya Janata Party (BJP) has been far less averse to nuclear weapons \citep{hymans_psychology_2006}. With the BJP leading the government for the first extended period of time, Prime Minister Atal Vajpayee authorized India's second nuclear weapons tests in 1998, despite security analysts widely believing that doing so unnecessarily inflamed tensions in South Asia \citep{bajpai_indian_2014}. Even as he authorized nuclear weapons testing, Vajpayee attempted to maintain India's reputation as a vocal proponent of disarmament. In a statement to parliament outlining nuclear doctrine after testing, Vajpayee argued that not only was India a responsible nuclear state, but that “our strengthened capability adds to our sense of responsibility.” He continued to argue that, \blockquote{the present decision and future actions will  continue to reflect a commitment to sensibilities and  obligations of an ancient civilization, a sense of responsibility and restraint, but a restraint born of the assurance of action, not of doubts or apprehension.}

In 2002, India would more explicitly incorporated “restraint” into its nuclear doctrine. Declaring a No First Use policy, Vajpayee said that India would refrain from using nuclear weapons unless first attacked by biological, chemical, or nuclear weapons. Today, there is widespread doubt regarding whether India would actually adhere to its NFU during a crisis with Pakistan \citep{clary_indias_2019, sundaram_india_2018}. In its 2014 election manifesto, the BJP said that India's NFU should be “revised and updated.” When asked by Reuters journalists what this revision entailed, an anonymous BJP source said to “read the manifesto” \citep{sanjeev_miglani_indias_2014}. After public backlash, Narendra Modi clarified that India would maintain its NFU, but since he became prime minister, Modi's government has continued to raise doubts that India will abide by its NFU. In 2019, after a terrorist attack in Indian-controlled Kashmir, India struck non-disputed Pakistani territory for the first time since 1971. Then-U.S. Secretary of State Mike \citet{pompeo_never_2023}
 claimed in his memoir that had the U.S. not been involved, the conflagration risked spiraling and leading to nuclear use. Several months after the attack Indian Defense Minister Rajnath Singh, standing at the site of India's second nuclear tests, announced that although India had previously “strictly adhered” to its NFU, “what happens in future depends on the circumstances” \citep{.sadf}. In May 2025, India and Pakistan engaged in a skirmish similar to that in 2019, during which India struck a Pakistani air base near the military headquarters responsible for overseeing its nuclear arsenal. India has also built out its intelligence gathering operations, improved the range of its ballistic missile system, and developed multiple targetable re-entry vehicles that would enable it to launch a counterforce attack against Pakistan that would disable its nuclear weapons facilities \citep{clary_indias_2019}.

Even as public officials have raised doubt about the legitimacy of India's NFU, they have clung to the language of India being a “responsible” nuclear state. In the same speech where he questioned that India would abide by its NFU, Singh said that “India attaining the status of a responsible nuclear nation is a matter of national pride.” After public outcry over the BJP's election manifesto, Modi called India's NFU “a reflection of cultural inheritance” \citep{busvine_indias_2014}. The contradictory nature of these statements suggests that although Indian elites may be disillusioned with the utility of India's NFU, they still believe that it is a useful signaling device either to domestic political suspicious of an aggressive nuclear posture or adversaries skeptical of Indian nuclear use.

\section{Research Design}
To test my hypotheses, I conduct a conjoint survey experiment modeled off of those created by Sagan and Valentino. In the experiment, respondents read a scenario describing a hypothetical Indian counterforce attack against Pakistan. The situation described is consistent with India's Cold Start doctrine and inspired by what security scholars consider to be the most likely scenario in which India would first use nuclear weapons \citep{kapur_ten_2008, clary_indias_2019}. Implemented after the Kargil War, Cold Start is military doctrine enabling India to quickly mobilize along the border with Pakistan following a state-sponsored terrorist attack. Within days of mobilization, India can push across the borders with Punjab and Kashmir either to claim disputed territory or force a favorable compromise. 

Although India publicly maintains a no first use policy, Christopher Clary and Vipin Narang argue that during such a ground war, India may contemplate launching a nuclear attack. Because India's military is conventionally superior to Pakistan's, the latter could consider using a small-yield tactical nuclear weapon on the battlefield to stop an Indian incursion. However, if Pakistan is incentivized to use nuclear weapons, India is encouraged to take first mover advantage and do the same. Fearing an eventual strike against one of its cities, Clary and Narang believe that India could launch a counter force attack, where it tries to eliminate Pakistan's strategic nuclear weapons. With the risk of a large-scale Pakistani attack marginalized, India could then conduct the ground war on its own terms.

\textit{A priori}, it is difficult to know how successful such an attack would actually be, given uncertainty over where Pakistan's nuclear weapons facilities are, how effective its air defenses would be, and the ordinary risks entailed by any large-scale military operation. In the article, I note that it is unclear whether all of Pakistan's strategic nuclear weapons have been destroyed, but that Indian government officials speaking on the condition of anonymity believe that there is a low chance of retaliation. I do so as to create a more difficult test for the norm against nuclear use while not misleading respondents into believing that a counterforce attack would be guaranteed to work. In a similar vein, although it is difficult to gauge how many civilians would die in a counterforce attack, I note that the World Health Organization warns that it could reach into the hundreds of thousands. Again, I do so as to not mislead respondents into believing that a counter force attack could be committed with little harm to civilians, even if the extent of that harm is ambiguous.

The story is structured like a mock news article so that key pieces of information can be highlighted in the headline, dek, and pull quotes.\footnote{All experimental materials can be found in  \autoref{app:vignettes}.} Specifically, I emphasize information that could influence people's understanding of the strategic utility of the attack and whether it was morally justifiable—including the target of the attack, that it was launched during an ongoing ground war, the number of civilian causalities, and the risk of retaliation. Across treatment and control groups, the article is identical, except that people assigned to learn about India's NFU read an additional sentence saying that “the nuclear attack violates long-standing Indian policy that it would never use nuclear weapons unless first attacked with weapons of mass destruction.” A pull quote in the middle of the article draws further attention to this fact.

After reading the vignette, respondents are randomly assigned to either reading or not reading three arguments against nuclear use. The arguments are written to invoke India's personal, positive identity and its negative, social identity vis-a-vis nuclear use. The first two directly mention Gandhi, Nehru, and India's legacy of non-violence. In addition to calling on these historical figures, the second and third juxtapose India with other states, saying that the use of nuclear weapons violates India's unique history as an advocate for global disarmament. The second explicitly mentions the U.S. and U.K. Mentioning the only state to have used nuclear weapons and the one that colonized the subcontinent may serve to bolster nationalist claims that India has a unique moral authority. In both the first and second argument, there are appeals in the first-person plural. Evoking a conception of “we” is important because norms are inter-subjective. It is not sufficient for one person to hold a reservation about a given course of action, but it must be held among a group of people. The third argument does not include specific mentions of Indian historical figures or “we-ness”, but refers broadly to India as having been a responsible state and upholding moral leadership. 

To measure how strongly the public supports nuclear weapons use, I include six different survey questions. In keeping with existing survey experiments, I ask respondents how strongly they approve of the nuclear attack as described in the vignette and how strongly they would prefer it to an alternative (here, continuing the ground war with Pakistan). Approval is the most direct way to measure whether someone supports nuclear use. Approval is also relevant to elected politicians fearful of domestic backlash and keen on limiting public opposition to their actions. Asking people whether they would prefer nuclear use to another military action is valuable because it reveals how strongly they have internalized just war theory. Whereas the nuclear attack described could not be defended under just war theory because of its failure to discriminate between combatants and civilians, as well as its lack of proportionality, the ground war would be defensible depending on the style and extent of combat, which is not described in the vignette. 

In addition to asking for people's approval and preferences, I ask two questions regarding what material action they would take to punish or not punish the sitting government. Specifically, would they be more or less likely to vote for the ruling party, and how likely would they be to protest nuclear use? Similar questions have yet to be asked in extant survey experiments, but offer a more direct test of whether politicians would face reprimand for nuclear use than simply asking for public approval. Though approval is a useful measure of overall public support for a politician, it is only relevant to an election-seeking politician who believes that it will affect how people vote. If a given policy decision is widely unpopular but low salience, then it will have little effect on voting patterns. By contrast, if the reaction to an action is polarized but high salience, then it could either incentivize or disincentivize a politician from supporting nuclear use, depending on the distribution of voters who support it. 

Lastly, I include measures of how moral respondents believe the attack to be. One of the questions asks directly whether respondents believe that the nuclear attack was ethical. Another asks whether they think it violates \textit{ahimsa}. These questions are useful in so far as they may reveal consequentialist reasoning. If a respondent believes that an attack is not moral but approves of it nonetheless, then that could imply they are willing to justify the attack based on the military and strategic benefits that it produces.

All outcome variables are measured on seven point Likert scales, with respondents being given the option to select “don't know.” If I were to find evidence to substantiate my hypotheses, I could observe two different, potentially compatible, patterns in people's responses. On one hand, support for nuclear use could decline across the seven point Likert scale. In this scenario, people's support for nuclear use changes with little difference in the number of survey takers reporting that they don't know how to respond. In another scenario, the incidence of people reporting that they don't know how to respond could increase while the number of people supporting the attack decreases. In this case, identity based arguments convert people into fence sitters. Instead of being in favor of military action, they are unsure of how to respond. To account for this possibility, I run all of my analyses twice. In one case, I exclude all NA observations; in the other, I impute values for people who say “don't know,” placing them in the middle of the Likert scale for each of the my outcomes.

Using inverse covariant weighting (ICW), I create an index that represents respondents' cumulative opposition to nuclear use. To create this index, I use the public's approval of the attack, whether it would affect how they vote, and whether they would publicly protest it. I group these factors together because they are indicative of how strongly a politician would be sanctioned for nuclear use. Creating such an index is useful to my study both practically and theoretically. ICW is useful practically because it enhances statistical power and reduces the number of tests that I run, minimizing the multiple comparison problem. Theoretically, ICW is useful because an election-seeking politician will weigh several aspects of public opposition to a nuclear attack when deciding whether to conduct one—not strictly, for example, public disapproval of the attack or whether people would protest it. As a result, though ICW does not provide an easily interpretable outcome variable, the index is holistic and provides insights into the different possible sources of opposition to nuclear use. As such, I estimate the following OLS model as the primary means to evaluate hypotheses two through four:

\begin{equation}
Y_i=\alpha+\beta_1Identity_i+\beta_2NFU_i+\beta_3(Identity_i*NFU_i)+\mathbf{\delta'X_i} +\epsilon_{i},
\end{equation}

\noindent
where $i$ indexes respondents, $Y_i$ represents composite attitudes toward nuclear use, $\beta_1$ is the fixed effect of reading identity-based arguments, $\beta_2$ is the fixed effect of telling respondents about India's NFU, $\beta_3$ is the fixed effect of the interaction between both treatments, and $X_i$ is a vector of individual-level covariates, including gender, religion, caste, age, income, and party. $Y_i$ is coded such that the higher values reflect stronger opposition to nuclear use.

Although the above model is useful for gauging whether my treatments have an effect on public opposition to nuclear use in the manner that I expect, it has several limitations. Most notably, it cannot be used to evaluate my first hypothesis. To determine whether the public has a categorical aversion to nuclear use, we must look at the sheer magnitude of support in the control group. I do so by collapsing public approval for nuclear use into three categories: disapproval, indifference, and approval. I reduce approval of the attack into three categories instead of two because indifference does not complement arguments either that there is a taboo against nuclear use or that people are particularly receptive to it. As a result, to generate a descriptively accurate assessment of public opposition to nuclear use, I do not collapse indifference into either approval or disapproval of the attack.

Second, even though ICW is a useful way to create a composite score for support of nuclear use, it weights different measures in a manner that's theoretically arbitrary. To be sure, factor analysis or another means of creating a composite score for public opposition to nuclear use would struggle from the same problem. Because different leaders would not weigh different sources of public opposition to nuclear use in the same manner, any particular weighting system that researchers impose would not universally apply. That is to say nothing of the fact that because ICW and factor analysis do not produce easily interpretable results, they would not be useful for political elites deciding whether to use nuclear weapons. 

I account for these problems in two ways. First, to test the robustness of my findings, I implement principal component analysis (PCA). Although this approach does not resolve issues regarding the interpretability of my outcome variable, it will show whether my treatments are still statistically significant using a different weighting scheme, reducing concerns that any statistically significant finding is attributable to one configuration of my composite outcome. Second, I estimate a series of logistic regression models, in which I dichotomize every outcome variable such that indifference to and disapproval of the attack are grouped together. Logistic regression is useful because ICW and PCA are “black boxes.” That is to say, it is not clear which measures of my outcome are being up- or down-weighted. By breaking apart respondents' answers across each of my outcomes, I evaluate how my treatments affect each type of opposition distinct from one another. Although this approach raises the problem of multiple comparisons, it reduces concerns regarding the theoretical arbitrariness of using any given weighting scheme. As such, in addition to my first model, I also estimate the following equation:

\begin{equation}
\log(\frac{\Pr(Y_i=Support)}{\Pr(Y_i\neq Support)})=\alpha+\beta_{1}Arguments_i+\beta_{2}NFU_{i}+\beta_{3}(Arguments_i*NFU_i)+\mathbf{\delta'X_i}
\end{equation}

\noindent
where $Y_i$ whether $i$ supports nuclear use across each measure of my outcome that I named above. All other parameters represent the same as they did in the previous equation. In text, I report the results for equations (1) and (2). In the Appendix, I include further robustness tests.

To conduct my experiment, I use Amazon Mechanical Turk to recruit two thousand participants to fill out a survey on Qualtrics. I select MTurk primarily for its cost efficiency compared to other survey platforms. Being cheap to recruit respondents, MTurk enables me to gain enough statistical power to adequately test the hypotheses that I proposed above. That being said, though highly correlated with nationally representative surveys, MTurk's results are not as representative as either probability sampling or online survey panels \citep{berinsky_evaluating_2012}. In India, MTurk's samples are disproportionately young, male, upper-caste, and wealthy \citep{boas_recruiting_2020}. To account for demographic differences, I weight my survey results using entropy balancing. 

That being said, because my survey is conducted in English, it cannot reasonably be construed as representative even after weighting. Because English is most widely spoken among people who are well educated, my findings will likely only apply narrowly to upper-class, elite opinion. In terms of the treatment effect, it is plausible that identity-based appeals could be salient in elite opinion but not in the general population. Given how well educated English-speaking people are on balance, then it is reasonable to believe that they know more about Indian history and foreign policy than the median voter. As a result, identity-based appeals could be more salient in my sample than they would be in the population writ large.

\section{Results}

\newpage
\begin{landscape}
% Table created by stargazer v.5.2.3 by Marek Hlavac, Social Policy Institute. E-mail: marek.hlavac at gmail.com
% Date and time: Tue, Mar 10, 2026 - 21:42:38
\begin{table}[!htbp] \centering 
  \caption{Logistic Regression Models} 
  \label{} 
\begin{tabular}{@{\extracolsep{5pt}}lccccccc} 
\\[-1.8ex]\hline 
\hline \\[-1.8ex] 
 & \multicolumn{7}{c}{\textit{Dependent variable:}} \\ 
\cline{2-8} 
\\[-1.8ex] & Approve & Prefer & Vote & Protest & National Identity & Ahimsa & Ethical \\ 
\\[-1.8ex] & (1) & (2) & (3) & (4) & (5) & (6) & (7)\\ 
\hline \\[-1.8ex] 
 NFU & 0.070 & 0.243 & 0.071 & $-$0.248 & $-$0.220 & $-$0.561$^{*}$ & 0.0004 \\ 
  & (0.307) & (0.301) & (0.303) & (0.304) & (0.311) & (0.311) & (0.298) \\ 
  & & & & & & & \\ 
 Identity Arguments & 0.328 & 0.004 & 0.289 & $-$0.149 & $-$0.425 & $-$0.417 & 0.016 \\ 
  & (0.327) & (0.310) & (0.318) & (0.314) & (0.322) & (0.324) & (0.312) \\ 
  & & & & & & & \\ 
 NFU x Identity & $-$0.197 & 0.176 & $-$0.293 & 0.181 & 0.463 & 0.344 & 0.634 \\ 
  & (0.449) & (0.433) & (0.436) & (0.433) & (0.441) & (0.438) & (0.437) \\ 
  & & & & & & & \\ 
 Constant & 0.317 & $-$0.356 & $-$0.501 & 0.602 & 0.564 & 1.321$^{**}$ & 0.465 \\ 
  & (0.634) & (0.613) & (0.617) & (0.615) & (0.629) & (0.638) & (0.636) \\ 
  & & & & & & & \\ 
\hline \\[-1.8ex] 
Observations & 419 & 419 & 419 & 419 & 419 & 419 & 419 \\ 
Log Likelihood & $-$253.003 & $-$265.449 & $-$262.205 & $-$264.457 & $-$258.771 & $-$261.633 & $-$260.801 \\ 
Akaike Inf. Crit. & 576.006 & 600.899 & 594.410 & 598.913 & 587.541 & 593.266 & 591.602 \\ 
\hline 
\hline \\[-1.8ex] 
\textit{Note:}  & \multicolumn{7}{r}{Fixed effects for caste, education, income, party, religion, gender, and age included but not shown} \\ 
 & \multicolumn{7}{r}{* p<0.1; ** p<0.05; *** p<0.01} \\ 
\end{tabular} 
\end{table} 

\end{landscape}

\newpage
\begin{landscape}
% Table created by stargazer v.5.2.3 by Marek Hlavac, Social Policy Institute. E-mail: marek.hlavac at gmail.com
% Date and time: Wed, Mar 11, 2026 - 12:17:15
\begin{table}[!htbp] \centering 
  \caption{OLS Models} 
  \label{} 
\begin{tabular}{@{\extracolsep{5pt}}lcccc} 
\\[-1.8ex]\hline 
\hline \\[-1.8ex] 
 & \multicolumn{4}{c}{\textit{Dependent variable:}} \\ 
\cline{2-5} 
\\[-1.8ex] & ICW & PCA & attack\_index\_likert\_pca & attack\_index\_na\_pca \\ 
 & \multicolumn{2}{c}{NA Imputation} & \multicolumn{2}{c}{Midpoint Imputation} \\ 
\\[-1.8ex] & (1) & (2) & (3) & (4)\\ 
\hline \\[-1.8ex] 
 NFU & $-$0.083 & $-$0.083 & 0.142 & 0.063 \\ 
  & (0.056) & (0.100) & (0.149) & (0.275) \\ 
  & & & & \\ 
 Identity Arguments & $-$0.031 & $-$0.020 & $-$0.105 & $-$0.378 \\ 
  & (0.058) & (0.109) & (0.155) & (0.301) \\ 
  & & & & \\ 
 NFU x Identity & 0.090 & 0.013 & $-$0.008 & 0.357 \\ 
  & (0.080) & (0.150) & (0.214) & (0.412) \\ 
  & & & & \\ 
 Constant & $-$0.057 & 0.034 & 0.458 & $-$0.455 \\ 
  & (0.116) & (0.220) & (0.307) & (0.605) \\ 
  & & & & \\ 
\hline \\[-1.8ex] 
Observations & 419 & 161 & 419 & 161 \\ 
R$^{2}$ & 0.051 & 0.128 & 0.085 & 0.198 \\ 
Adjusted R$^{2}$ & $-$0.033 & $-$0.107 & 0.004 & $-$0.018 \\ 
Residual Std. Error & 0.395 (df = 384) & 0.402 (df = 126) & 1.051 (df = 384) & 1.104 (df = 126) \\ 
F Statistic & 0.607 (df = 34; 384) & 0.544 (df = 34; 126) & 1.046 (df = 34; 384) & 0.917 (df = 34; 126) \\ 
\hline 
\hline \\[-1.8ex] 
\textit{Note:}  & \multicolumn{4}{r}{Fixed effects for caste, education, income, party, religion, gender, and age included but not shown} \\ 
 & \multicolumn{4}{r}{* p<0.1; ** p<0.05; *** p<0.01} \\ 
\end{tabular} 
\end{table}
\end{landscape}
\newpage

\bibliography{references}

\appendix

\section{Tables}

\section{Experimental Materials}
\label{app:vignettes}
\subsection*{Control Article}
\includepdf[pages=-]{controlpdf.pdf}
\subsection*{Treatment Article}
\includepdf[pages=-]{nfupdf.pdf}
\subsection*{Arguments}

\textbf{Reaction 1:} The government’ s use of nuclear weapons was morally repugnant and contradicts Mahatma Gandhi’s, Jawaharlal Nehru’s, and this entire country's legacy of non-violence. It violates everything that we know about ourselves—it should never have happened.

\textbf{Reaction 2:} India for decades has led the fight for a nuclear-free world, criticizing the Americans, British, and whoever else for their cavalier attitude toward nuclear weapons. Gandhi and Nehru wanted a safer world, a more dignified world, where the annihilating power of nuclear weapons didn’t haunt anyone. This government’s decision to use the same weapons, however, has made India abandon that legacy. We are now no better than rest.

\textbf{Reaction 3:} For generations, India claimed moral leadership by rejecting the logic that security must rest on mass destruction. That claim carried weight precisely because others failed to uphold it. By resorting to nuclear violence, this government has dissolved the distinction between India’s vision of global responsibility and the reckless behavior it once condemned.

\section{Questionnaire}

\end{document}